% Appendix C: numerical checks (readability layer, decision D10). The remarks of the text keep their
% numbers and point here. No numbered objects other than the appendix section itself.
\section{Numerical checks}\label{app:numerics}

This appendix describes the numerical checks mentioned in the text. No proof uses them. Each paragraph names the place in the text that refers to it. The programs and their outputs are in the folder \texttt{numerics} next to the note, and \texttt{numerics/README.md} describes them.

\paragraph{The averaging of Part~IV (Remark~\ref{rem:avg-numerical}).}
The algebra of Sections~\ref{sec:averaging} and~\ref{sec:adapted} was checked numerically for $G=SU(2)$ and $L=3$: in $d=2$ on small tori, with $k=1$ and $k=2$, and in $d=4$ on a single coarse bond, where the families have up to $24$ contours. The following hold to rounding error:
\begin{itemize}
\item for real families of $2$ to $24$ elements, Lemma~\ref{lem:group-average}(a), (b) with real $u,v$, (c), and $\sum_i\omega_i=I$ with $|\omega_i-\frac1nI|$ proportional to the diameter;
\item \eqref{eq:Y2-law} for real gauge transformations that are not small;
\item Lemma~\ref{lem:reversal}, Lemma~\ref{lem:one-step-basic}(b) for real $g$, and Lemma~\ref{lem:one-step-basic}(d) for all $384$ elements of $B_4$;
\item Theorem~\ref{thm:correspondence}(e),(f) and Proposition~\ref{prop:linearized-averages}(a),(b),(f);
\item at $U_0=1$, Lemma~\ref{lem:one-step-expansion}(d),(f), with the coefficient $L^{1-d}$, and Proposition~\ref{prop:linearized-averages}(c).
\end{itemize}
Two controls fail, as they should. With a single contour order per family, the covariance under the permutations of the axes fails. With the mean of the logarithms of the contour variables in place of $\Av$, \eqref{eq:Y2-law} and $E=0$ fail. Not tested numerically: the complex equivariance in Lemma~\ref{lem:group-average}(b), the formula for the weights in (d), the Lipschitz bound (e), complex $g$ in Lemma~\ref{lem:one-step-basic}(b), and Lemma~\ref{lem:avg-transport}(e),(f).

\paragraph{The constant of the one-step Poincar\'e inequality (Section~\ref{sec:covariance}).}
Finite samples of the Bloch spectrum relevant to the constant $C_{\rm P}$ of Lemma~\ref{lem:poincare-step} were computed numerically at $W=1$. By Bloch's theorem over $\T'$ the form $\|\dl B\|^2+\|\mathfrak G_1B\|^2+\|q^{\rm pl}B\|^2$ reduces to finite matrices, one for each momentum. For $d=4$ and $L=3$ its smallest eigenvalue lies in $[0.030,0.035]$ over the momenta tested, with the minimum among those samples about $0.030$ at $(\pi,\pi,\pi,\pi)$; the reciprocal of this sampled minimum is approximately $33$. These finite samples do not certify the minimum over all Bloch momenta or a uniform Poincar\'e constant of that size. The existence of a constant uniform in the size of the torus follows from the exact proof of Lemma~\ref{lem:poincare-step}, independently of these numerical values. In $d=2$ the corresponding sampled smallest eigenvalues are about $0.12$ for $L=3$ and $0.08$ for $L=5$.

\paragraph{Gauge fixing and the one-loop term (Remark~\ref{rem:one-loop-intrinsic}).}
For $d=2$, $G=SU(2)$, $L=3$ and four coarse fields, Laplace's method with the exponential gauge fixing over all bonds was compared with Laplace's method in a tree gauge, which has no Faddeev--Popov determinant either. The two differ by the constant $\frac{n_\lambda}2\log2$, where $n_\lambda=3|\T\setminus\T'|$ is the number of gauge directions, to $2.5\cdot10^{-10}$; this is the exact prediction for a gauge fixing without Faddeev--Popov determinant. The constant comes from the normalizations of the program, which omits the factor $z^{-1}$ of $\mathrm{gf}_W$ and uses the coordinates $h=e^{\ii H}$, $H=\frac12\sum_ax_a\tau_a$ with the Pauli matrices $\tau_a$, so that $\operatorname{tr}H^2=\frac12|x|^2$. In these coordinates the exponential weight at the coupling of the action is a Gaussian of variance $2g^2$ in each eliminated gauge direction, which contributes $\frac12\log2$ relative to the tree gauge. Two exact representations of one integral have the same Laplace asymptotics. The program is \texttt{symavg.py} in the folder \texttt{numerics} next to the note.

\paragraph{The one-loop term in the Landau chart (Remark~\ref{rem:ol-numerical}).}
We tested Theorem~\ref{thm:ol-chart} numerically for $d=2$, $G=SU(2)$ and $k=1$. This is outside the hypothesis $d=4$ of the note, but the identity is algebraic. The first test, with the program \texttt{chart.py}, was made for an earlier, simpler averaging, with blocks that have the coarse point at a corner and a single coordinate-ordered contour per point, on three lattices: $L=2$ on coarse tori of $2\times2$ and $3\times3$ points, and $L=3$ on a coarse torus of $2\times2$ points. The second test is described at the end of the remark. Each coarse field $V$ was the average $U^{(1)}$ of a random fine configuration. For each $V$:
\begin{itemize}
\item A critical point $U_0$ of $S$ on the fibre $\{U^{(1)}=V\}$ was computed by a damped Newton method for the Lagrangian, which accepts only steps that decrease $S$. It was computed in a tree gauge, where $U=1$ on the bonds of a spanning tree of each block, formed by coordinate-ordered contours from the coarse point.
\item The left side was computed by Laplace's method in the same gauge. The tree gauge is a computational device of the programs and plays no role in the note: by the invariance of the Haar measure, the variables on the tree bonds can be fixed one at a time, so the Faddeev--Popov identity for this gauge has Jacobian $1$. The left side is $-\frac12\log\det$ of the Hessian of the Lagrangian on the tangent space of the fibre, minus $\frac12\log\det$ of the Gram matrix of the differential of the constraint.
\item The right side was computed from the gauge-adapted averaging $\mathcal Q$ of the program, relative to $U_0$. In it, $q_2$ includes the multiplier term, $R$ is the projection onto $\cW=D^TD\,\cN(Q'_1)$, and the ghost is $\frac12\log\det(PX^2P|_{\cN(Q'_1)})$, as in the proof of Corollary~\ref{cor:ol-ghost}.
\end{itemize}
For the averaging of the first test the tree slice is the axial slice. On it, $\mathcal Q$ coincides with the constraint $U^{(1)}V^{-1}$, as in Theorem~\ref{thm:correspondence}(d), so the two sides share the constraint and the multiplier. What the first test checks is the passage from the axial slice to the Landau slice: the Faddeev--Popov factors, the power of the ghost, the term $\bQ_1\bG_*\bQ_1^T$, and the multiplier term in the directions off the axial slice. It does not test Theorem~\ref{thm:chart}, which only justifies the chart.

First derivatives were taken by complex steps on the holomorphic extensions, which give them to machine precision. Second derivatives combine a complex step with a central difference, with accuracy about $10^{-10}$.

Results:
\begin{itemize}
\item On each lattice the difference of the two sides was independent of $V$ to within $5\cdot10^{-10}$. Meanwhile each side varied by up to $2.1$ between the coarse fields: seven of them for $L=2$ on the $2\times2$ torus, and two on each of the other lattices. The difference is in fact exactly $-\frac d2\dim\mathfrak g\,|\T|\log L$, i.e.\ $-12\log2$, $-27\log2$ and $-12\log3$. The program imposes the coarse condition in unnormalized Lie coordinates, in which the block $\Theta_0$ of $\mathcal O_0$ in Proposition~\ref{prop:ol-FP} is $L^{-d/2}Q'_1$ instead of $Q'_1$. So the test checks an exact identity, not only the independence of $V$.
\item The right side did not depend on $a$.
\item The identities used in the proof hold to within $1.2\cdot10^{-9}$: $\nabla\Phi(0)=\bQ_1^T\omega_*$ (Proposition~\ref{prop:euler-lagrange}), and $q_2D\lambda=0$ and $\bQ_1D\lambda=0$ for $\lambda\in\cN(Q'_1)$.
\item The identity of Corollary~\ref{cor:ol-ghost} holds to $10^{-13}$, at a random background and for an averaging with $Q'Q'^T=I$.
\end{itemize}
The test can detect errors. Wrong variants of the formula give differences that vary with $V$ by $1.9\cdot10^{-2}$ to $1.7$: without the ghost, with its square, without the multiplier term in $q_2$, and without $-\frac12\log\det(\bQ_1\bG_*\bQ_1^T)$. The programs and their outputs are in the folder \texttt{numerics} next to the note. The program \texttt{symavg.py} there repeats the test for Ba\l aban's symmetric averaging in four dimensions, with centred blocks and averaged contour variables, transcribed to $d=2$; this is the averaging of Part~IV, with the averaged axial gauge and the gauge-adapted averaging of Section~\ref{sec:adapted}. It finds the same identity with the same constant: for $L=3$ on a coarse torus of $2\times2$ points and four coarse fields, between which each side varies by $0.64$, the difference of the two sides is $-12\log3$ to within $5\cdot10^{-11}$. Here the tree slice and the averaged axial slice differ, so this test checks the identity as a whole. With a single ordered contour in place of the averaged contour variables in the gauge, the right side is the same to $5\cdot10^{-11}$, as it must be for any admissible chart. With logarithmic means of the contour variables in place of the group average, for which the exact covariance $E=0$ fails, the difference misses the prediction by $2\cdot10^{-4}$ to $7\cdot10^{-4}$.

\paragraph{Numerical evaluation of $f$ (Remark~\ref{rem:ol-f-numerical}).}
For $G=SU(2)$ in the fundamental representation, $\mathsf b=\frac{11}{6\pi^2}=0.185756$. We evaluated $f(\mathrm m)$ numerically, with a real background along a Cartan generator, $A_2(x)=\zeta\,e^{\ii p_1x_1}+\bar\zeta\,e^{-\ii p_1x_1}$ and $A_\mu=0$ for $\mu\ne2$. The vacuum polarization is the coefficient of $|\zeta|^2$, per site.
\begin{itemize}
\item The vertices of the Wilson term come from the exact Taylor coefficients of the plaquette summand in the chart $e^{\ii h}e^{\ii A}$, up to order $h^2A^2$. These were computed by Cauchy's formula on the holomorphic extension, to accuracy $10^{-10}$. The vertices of $D_UD_U^T$ and $D_U^TD_U$ come from the expansion of $R(e^{\ii A})$. In the component pairing the propagators are $2/(|\hat q|^2+\mathrm m^2)$ for the vector field and $1/(|\hat q|^2+\mathrm m^2)$ for the scalar field.
\item The $p_1^2$-coefficient was taken as a second difference in $p_1$ at the level of the integrand, with steps $10^{-3}$ and $4\cdot10^{-3}$. Richardson extrapolation removed the bias of order (step)$^2$, which was up to $4\cdot10^{-6}$ at the smallest $\mathrm m$.
\item The integral over $[-\pi,\pi]^4$ was computed by the trapezoidal rule with $N^4$ points, $N\mathrm m\ge14$. For smooth periodic integrands this rule converges exponentially. Comparisons of two values of $N$, made for three of the masses, put the remaining errors of the values below at less than $10^{-6}$.
\end{itemize}
The results are:
\begin{center}\small
\begin{tabular}{@{}lcccccccc@{}}
\hline
$\mathrm m$ & $1$ & $0.5$ & $0.354$ & $0.25$ & $0.177$ & $0.125$ & $0.088$ & $0.0625$ \\
$f(\mathrm m)$ & $0.50315$ & $0.67048$ & $0.74416$ & $0.81396$ & $0.88141$ & $0.94750$ & $1.01281$ & $1.07770$ \\
\hline
\end{tabular}
\end{center}
Three conclusions:
\begin{enumerate}[label=(\roman*)]
\item \emph{Ward identity.} $\Pi(0)$ vanishes to $10^{-9}$ for all $\mathrm m$. The longitudinal part $\Pi_{11}(p_1)$ at $p_1=0.5$ and $\mathrm m=1$ vanishes to $10^{-10}$, the accuracy of the Taylor coefficients. The tadpole and bubble terms are separately of order $0.3$. Since each of the two determinants is gauge invariant on its own, this checks the Wilson and gauge-fixing vertices against each other, and the ghost vertices separately. The relative weight of the ghost is checked by (ii) and (iii).
\item \emph{The logarithm and its sign.} Fits of $f(\mathrm m)=\mathsf b\log(1/\mathrm m)+c+\mathrm m^2(d_2\log(1/\mathrm m)+e_2)$, with or without terms of order $\mathrm m^4$, give $\mathsf b$ between $0.18564$ and $0.18579$, against $0.185756$. The fits that use only $\mathrm m\le\frac14$, and the fits with six parameters, give $0.18574$ to $0.18577$. The function $f$ increases as $\mathrm m\to0$. This is the sign of asymptotic freedom.
\item \emph{The constant.} With $\mathsf b$ fixed, the fits give $c$ between $0.563268$ and $0.563274$. Dashen and Gross, and Hasenfratz and Hasenfratz, computed the one-loop relation between the bare Wilson coupling and the $\overline{\rm MS}$ coupling \cite{DashenGross,HasenfratzLambda}. For $SU(N)$ it reads $\Lambda_{\overline{\rm MS}}/\Lambda_L=38.852704\,e^{-3\pi^2/(11N^2)}$, that is, $g_{\overline{\rm MS}}^{-2}(a^{-1})=g_0^{-2}-d_1$ with $d_1=0.16995599N-\frac1{8N}$, in the normalization $\beta=2N/g_0^2$. Combined with the $\overline{\rm MS}$ value of the continuum term of Lemma~\ref{lem:ol-continuum-integrand}, it predicts
\[
c=\frac{d_1}{T_\rho}+\frac{C_A}{48\pi^2T_\rho}=0.563267\qquad(SU(2)).
\]
The second term comes from the trace over the vector index, which is $d=4-2\varepsilon$ rather than $4$. So the pole of the orbital term leaves a finite remainder. The agreement is to $10^{-5}$.
\end{enumerate}
So the lattice integral of Lemma~\ref{lem:ol-scaling} reproduces a known lattice constant, independently of the arguments of this section. In particular, the normalization of Lemma~\ref{lem:ol-classical} is confirmed for $SU(2)$. The programs are in the folder \texttt{numerics} next to the note.

\paragraph{The integrand and the tensor (Remark~\ref{rem:ol-integrand-numerical}).}
Two further checks are in \texttt{numerics/y5c} (\texttt{numerics/README.md}, Section~5).
\begin{enumerate}[label=(\roman*)]
\item \emph{The spherical average, in exact arithmetic.} The program computes $\langle\mathcal I_0(\omega,0)\rangle$ in rational arithmetic, from the vertices of Lemma~\ref{lem:ol-naive-limit} and the moments $\langle\prod_\alpha\omega_\alpha^{2b_\alpha}\rangle=\prod_\alpha(2b_\alpha-1)!!/(4\cdot6\cdots(2n+2))$, $n=\sum_\alpha b_\alpha$, which contain \eqref{eq:ol-sphere-moments}. All $256$ components of the four-index tensor have the form $a\,\delta_{\mu\nu}\delta_{\kappa\lambda}+b\,(\delta_{\mu\kappa}\delta_{\nu\lambda}+\delta_{\mu\lambda}\delta_{\nu\kappa})$, with $(a,b)=(\frac13,-\frac16)\kappa_\rho$ from the scalar determinant, $(-\frac23,\frac13)\kappa_\rho$ from the orbital part of the vector determinant, $(4,-2)\kappa_\rho$ from the spin part, and $(\frac{11}3,-\frac{11}6)\kappa_\rho$ in total. So $a=-2b$ holds for each part separately, and the isotropic coefficient is $\frac{11}3\kappa_\rho$, as in the proof of Lemma~\ref{lem:ol-classical}.
\item \emph{The scaling of the lattice integrand.} For $G=SU(2)$, with the lattice vertices of Remark~\ref{rem:ol-f-numerical}, the program compares $\lambda^4\mathcal I(\lambda\omega)$ with $\mathcal I_0(\omega)$, both contracted with a random real unit polarization $c$ and a random unit direction of $p$, for $200$ random directions $\omega$ on the closed half-sphere $\{|q|^2+\mathrm m^2=1,\ \mathrm m\ge0\}$ and the scale factor $\lambda=2^{-j}$, $0\le j\le12$. The second derivative in $p$ is a second difference with steps $10^{-3}\lambda$ and $4\cdot10^{-3}\lambda$, with Richardson extrapolation. The largest difference over the directions is $4.3$ at $\lambda=1$ and $2.4\cdot10^{-7}$ at $\lambda=2^{-12}$; it decreases like $4\lambda^2$, faster than the bound $C\lambda$ of Lemma~\ref{lem:ol-continuum-integrand}. The vertices of this test come from the Taylor coefficients of the plaquette, not from the proof of Lemma~\ref{lem:ol-naive-limit}. So it checks the limits of the vertices at the level of the integrand, including the factor $2$ of the spin term and the weight of the ghost, independently of the fits of Remark~\ref{rem:ol-f-numerical}.
\end{enumerate}

\paragraph{Summary (Remark~\ref{rem:ol-remains}(d)).} The checks of Part~XI, with the accuracies reached, are the following.
\begin{itemize}
\item Theorem~\ref{thm:ol-chart} was tested on small two-dimensional lattices, where it holds to $5\cdot10^{-10}$ (Remark~\ref{rem:ol-numerical}). This first test, made for a simpler averaging, checks the passage from the axial slice to the Landau slice.
\item Ba\l aban's symmetric averaging of Part~IV was transcribed to $d=2$, with $L=3$ (\texttt{numerics/README.md}, Section~3). Theorem~\ref{thm:ol-chart} holds on the averaged axial slice with the same exact constant $-12\log3$, to $5\cdot10^{-11}$. Laplace's method with the exponential gauge fixing of Lemma~\ref{lem:exp-gauge} differs from Laplace's method in a tree gauge by the constant $\frac{n_\lambda}2\log2$ predicted by the absence of a Faddeev--Popov determinant, to $2.5\cdot10^{-10}$ (Remark~\ref{rem:one-loop-intrinsic}). The one-loop term is invariant under a rotation, a reflection and the transposition of the axes, to $2.2\cdot10^{-10}$, as Lemma~\ref{lem:ol-symmetry} predicts. The algebra of the averaging, including its covariance under all $384$ elements of $B_4$ and the exact covariance $E=0$, was also checked in $d=4$ on a single coarse bond; there is no Laplace computation in $d=4$.
\item $f(\mathrm m)$ was evaluated for $SU(2)$ (Remark~\ref{rem:ol-f-numerical}). It confirms the coefficient $\mathsf b$ to $2\cdot10^{-5}$ (to $2\cdot10^{-4}$ for every fit), the sign, transversality, and the known lattice constant of Dashen and Gross.
\item The spherical average of the full continuum tensor was computed in exact rational arithmetic, and the scaling of the lattice integrand to the continuum integrand was checked for $SU(2)$ (Remark~\ref{rem:ol-integrand-numerical}; \texttt{numerics/README.md}, Section~5). The first gives the coefficients $\frac{11}3$ and $-\frac{11}6$ (in units of $C_A/T_\rho$) of the two invariant tensors, with transversality for the scalar, orbital and spin parts separately; in the second the difference decreases like $\lambda^2$ in the scale factor $\lambda$.
\end{itemize}
